\section{Bối cảnh và động lực thực hiện}
\subsection{Tổng quan về học sâu}
%Vị trí, tầm quan trọng của deeplearning trong lĩnh vực Machine learning
Học máy (Machine Learning) và học sâu (Deep Learning) là hai lĩnh vực đã được nghiên cứu từ lâu, nhưng chỉ thực sự bùng nổ trong những năm gần đây nhờ sự gia tăng đáng kể của dữ liệu và sự phát triển của các nền tảng phần cứng như GPU, NPU và hệ thống tính toán hiệu năng cao (HPC)\cite{ibm_2024_neural,26_sharifani_2023_machine}. Trong đó, học sâu đã trở thành một hướng tiếp cận chủ đạo của học máy nhờ khả năng tự động trích xuất đặc trưng và xử lý lượng lớn dữ liệu hiệu quả, kể cả dữ liệu phi cấu trúc. Học sâu được phát triển dựa trên nền tảng mạng nơ-ron nhân tạo nhiều lớp, lấy cảm hứng từ cách xử lý thông tin của não bộ. Hiện nay, lĩnh vực này tiếp tục phát triển với nhiều hướng tiếp cận đa dạng và các phương pháp học hiện đại mở ra tiềm năng ứng dụng rộng lớn trong nhiều lĩnh vực khác nhau\cite{28_khoei_2023_deep}. Theo như nhận định của Geoffrey Hinton, một trong những người tiên phong trong lĩnh vực trí tuệ nhân tạo, “học sâu sẽ có thể làm tất cả mọi thứ”\cite{hao_2020_ai}.

%Strength, weekness
Một trong những ưu điểm nổi bật của học sâu là khả năng tự động trích xuất đặc trưng, qua đó giảm đáng kể sự phụ thuộc vào các kỹ thuật lựa chọn đặc trưng thủ công trong học máy truyền thống. Bên cạnh đó, học sâu thể hiện hiệu quả vượt trội trong việc xử lý và khai thác thông tin từ lượng lớn dữ liệu phi cấu trúc như hình ảnh, âm thanh và ngôn ngữ tự nhiên,làm giảm các sai số tổng quát hóa so với các mô hình học máy truyền thống\cite{14_10.1145/3776588}. Trong một số bài toán phức tạp, độ chính xác học sâu đạt được thậm chí tiệm cận hoặc vượt qua con người. Tuy nhiên, phương pháp này vẫn tồn tại nhiều hạn chế như yêu cầu lớn về dữ liệu huấn luyện, nguy cơ sai lệch dữ liệu (bias), khả năng giải thích của mô hình, hiện tượng quá khớp (overfitting), chi phí tính toán,\cite{27_alzubaidi_2021_review} cùng với các vấn đề đạo đức và công bằng trong ứng dụng\cite{28_khoei_2023_deep}. Do đó, các nghiên cứu hiện nay vẫn đang phát triển các phương pháp, khung hướng dẫn và mô hình nhằm nâng cao tính hiệu quả cũng như đảm bảo tính đạo đức trong việc triển khai học sâu. 

%Application of deep learning
Học sâu, với năng lực tính toán và biểu diễn mạnh mẽ, đã tạo ra ảnh hưởng sâu rộng trong nhiều lĩnh vực như y tế, tài chính, an ninh mạng, robot và tự động hóa. Các ứng dụng tiêu biểu có thể kể đến như các hệ thống phân tích ảnh y khoa phục vụ phát hiện ung thư và COVID-19, trợ lý ảo thông minh, và xe tự hành. Bên cạnh đó, việc tích hợp học sâu với các công nghệ khác như IoT và blockchain đã mở ra nhiều hướng tiếp cận mới nhằm giải quyết các bài toán phức tạp trong thực tiễn\cite{27_alzubaidi_2021_review}. Đặc biệt, trong lĩnh vực thị giác máy tính và nhận dạng hình ảnh, học sâu - điển hình là các mạng nơ-ron tích chập (Convolutional neural network - CNN) - đã trở thành phương pháp chủ đạo nhờ các đặc tính như kết nối thưa, chia sẻ tham số và khả năng học biểu diễn hiệu quả\cite{Goodfellow-et-al-2016}.
\subsection{Tổng quan về Tiny Deep Learning (TinyDL)}
%Bối cảnh, tầm quan trọng về sự xuất hiện của tinyml & tinydl trên embedded system
Theo cách tiếp cận truyền thống, các thuật toán học máy và học sâu thường được triển khai trên các hệ thống có tài nguyên tính toán mạnh hoặc nền tảng điện toán đám mây. Trong những năm gần đây, nhằm giảm chi phí phần cứng và khắc phục các hạn chế của đám mây như độ trễ và quyền riêng tư, xu hướng triển khai mô hình trực tiếp trên thiết bị biên ngày càng gia tăng\cite{29_10177729}. Sự phát triển của các kỹ thuật nén mô hình, sự phát triển cũng như chi phí thấp của kiến trúc phần cứng đã thúc đẩy sự hình thành của TinyML - một hướng tiếp cận tập trung vào việc triển khai các mô hình học máy trên phần cứng hạn chế\cite{1_warden_2020_tinyml}. Trong đó, Tiny Deep Learning (TinyDL) là một nhánh chuyên biệt, hướng đến việc tối ưu và triển khai các mô hình học sâu, điển hình như mạng CNN, trên các nền tảng tài nguyên thấp, đặc biệt là vi điều khiển (MCU)\cite{14_10.1145/3776588}.

%Strength, weekness, challenges, application
Nhờ khả năng cung cấp phản hồi thời gian thực, tiêu thụ năng lượng thấp và đảm bảo quyền riêng tư, bảo mật dữ liệu thông qua xử lý cục bộ trên thiết bị, Tiny Deep Learning (TinyDL) đang được ứng dụng rộng rãi trong nhiều lĩnh vực như y tế, công nghiệp, môi trường, nhận diện từ kích hoạt và đặc biệt là các bài toán thị giác máy tính sử dụng camera độ phân giải thấp. Tuy nhiên, bên cạnh các giới hạn phần cứng cố hữu, các mô hình Tiny Deep Learning (TinyDL) vẫn đối mặt với nhiều thách thức, bao gồm rủi ro về an ninh, sự thiếu hụt các công cụ phát triển có tính di động cao tương thích với các nền tảng phần cứng không đồng nhất, cũng như nhu cầu tích hợp các kỹ thuật tối ưu hóa mô hình và các phương pháp học hiện đại như học trực tuyến và học liên kết. \cite{14_10.1145/3776588}

\subsection{Tổng quan về hệ thống nhận diện khuôn mặt}
%Tổng quan về hệ thống nhận diện khuôn mặt
Nhận diện khuôn mặt (Face Recognition) là một trong những ứng dụng quan trọng và được nghiên cứu rộng rãi trong lĩnh vực thị giác máy tính. Công nghệ này cho phép hệ thống tự động phát hiện, trích xuất đặc trưng và nhận dạng danh tính con người từ dữ liệu hình ảnh hoặc video, qua đó mở ra nhiều ứng dụng thiết thực trong các lĩnh vực như an ninh - giám sát, tài chính - ngân hàng, giao thông thông minh và quản lý đô thị \cite{30_9145558,ibm_2025_biometric}. Nhờ các ưu điểm về tính tiện lợi, khả năng tương tác tự nhiên, khả năng mở rộng và mức độ tự động hóa cao, công nghệ nhận diện khuôn mặt đã phát triển nhanh chóng và được ứng dụng rộng rãi trong nhiều lĩnh vực, từ mở khóa thiết bị di động, chấm công tự động, thanh toán không tiếp xúc đến giám sát tại các không gian công cộng \cite{17_10128985}. Về mặt kỹ thuật, bài toán này thường được phân thành hai hướng chính: định danh khuôn mặt (Face Identification) với bài toán so khớp 1:N, và xác minh khuôn mặt (Face Verification) với bài toán 1:1 nhằm kiểm tra tính xác thực của một danh tính đã khai báo; trong đó, xác minh khuôn mặt đặc biệt quan trọng trong các hệ thống tài chính ngân hàng với bài toán xác thực sinh trắc học yêu cầu độ chính xác cao và khả năng chống giả mạo. \cite{2_electronics9081188,techcombank_2025_xc}

Với sự phát triển mạnh mẽ của TinyML/TinyDL, các hệ thống nhận diện khuôn mặt ngày càng được triển khai tại biên nhằm giảm phụ thuộc vào điện toán đám mây và tăng cường bảo mật dữ liệu. Các mô hình nhẹ phục vụ bài toán nhận diện khuôn mặt cũng như nhận diện đa mô thức tiếp tục là hướng nghiên cứu được quan tâm. Tuy nhiên, việc triển khai trên các nền tảng vi điều khiển vẫn gặp nhiều thách thức do hạn chế về tài nguyên, đặc biệt là bài toán cân bằng giữa kích thước mô hình và độ chính xác, cùng với yêu cầu tích hợp các cơ chế chống giả mạo nhằm đáp ứng các tiêu chuẩn ứng dụng trong thực tế \cite{35_10457480,36_11134957}. Do đó, đặt ra nhu cầu nghiên cứu và hiện thực hóa một mô hình học sâu xác minh khuôn mặt có kích thước gọn nhẹ, phù hợp triển khai trên vi điều khiển, đồng thời đáp ứng yêu cầu trong cả thực tiễn và học thuật.

\section{Giới thiệu về đề tài}
\subsection{Tổng quan về đề tài}
Đề tài “Hệ thống xác thực danh tính khuôn mặt trên đa nền tảng” tập trung nghiên cứu và phát triển một giải pháp xác thực danh tính dựa trên khuôn mặt, có khả năng vận hành trên nhiều nền tảng phần cứng, bao gồm hệ thống vi điều khiển với tài nguyên hạn chế. Hệ thống sử dụng camera để thu nhận hình ảnh trực tiếp dưới dạng luồng dữ liệu, sau đó dữ liệu hình ảnh được xử lý tuần tự qua các mô-đun chức năng chính: phát hiện khuôn mặt (Face Detection), trích xuất đặc trưng (Feature Extraction) và xác minh khuôn mặt (Face Verification). Các mô-đun này phối hợp chặt chẽ nhằm đảm bảo quá trình xác thực diễn ra chính xác và nhanh chóng, đồng thời đáp ứng yêu cầu hiệu năng tính toán gọn nhẹ, phù hợp với các ứng dụng thời gian thực. Đề tài không chỉ tập trung vào việc thiết kế mô hình học sâu, mà còn chú trọng đến khả năng triển khai hệ thống trên đa dạng nền tảng phần cứng, đánh giá hiệu năng và tính ổn định trong các điều kiện vận hành thực tế.

\subsection{Lợi ích và sự cần thiết của đề tài}
Hệ thống xác thực danh tính sử dụng khuôn mặt mang lại nhiều lợi ích thiết thực cho người sử dụng cũng như cho các tổ chức, doanh nghiệp trong quá trình quản lý truy cập và đảm bảo an ninh. Trước hết, việc tự động hóa quá trình nhận diện và xác minh danh tính giúp loại bỏ hoàn toàn các thao tác thủ công dựa trên mật khẩu, thẻ từ hay mã PIN, từ đó giảm thiểu sai sót do con người, hạn chế nguy cơ lộ lọt thông tin và tiết kiệm đáng kể thời gian xác thực\cite{ibm_2025_biometric}. Nhờ khả năng nhận diện khuôn mặt theo thời gian thực, hệ thống cho phép xác thực nhanh chóng, không tiếp xúc, nâng cao tính tiện lợi cho người dùng, đặc biệt trong các môi trường yêu cầu lưu lượng truy cập lớn như tòa nhà thông minh, văn phòng, trường học, khu dân cư và hệ thống nhà ở thông minh (Smart Home).

Việc nghiên cứu và phát triển hệ thống xác thực khuôn mặt cũng mang tính cấp thiết trong bối cảnh chuyển đổi số cũng như sự phát triển của các hệ thống IoT đang diễn ra mạnh mẽ trên phạm vi toàn cầu. Xác thực sinh trắc học dần trở thành tiêu chuẩn trong các lĩnh vực như ngân hàng số, thương mại điện tử, quản lý hành chính công và an ninh đô thị. Tại Việt Nam, việc từng bước triển khai xác thực sinh trắc học trong các giao dịch tài chính, định danh điện tử và hệ thống quản lý thông minh cho thấy nhu cầu cấp bách về các giải pháp xác thực khuôn mặt có độ chính xác và độ tin cậy cao. Đặc biệt, đối với các hệ thống quản lý truy cập như Smart Home, một lỗ hổng trong quá trình xác minh danh tính có thể dẫn đến nguy cơ xâm nhập trái phép, gây mất an toàn cho tài sản và con người. Vì vậy, việc nghiên cứu và hoàn thiện hệ thống xác minh khuôn mặt là hết sức cần thiết, đặc biệt là trên các thiết bị nhúng và vi điều khiển.

Cuối cùng, xét từ góc độ học thuật và đào tạo, đề tài tạo điều kiện cho sinh viên ngành Kỹ thuật Máy tính tiếp cận quy trình xây dựng một hệ thống ứng dụng AI hoàn chỉnh, từ thu thập dữ liệu, tiền xử lý, huấn luyện mô hình, tối ưu thuật toán đến triển khai trên thiết bị vi điều khiển với tài nguyên giới hạn. Thông qua đó, sinh viên không chỉ được củng cố kiến thức chuyên môn về xử lý ảnh, học máy, học sâu và hệ thống nhúng, mà còn được trang bị kỹ năng thực tiễn, đáp ứng nhu cầu nhân lực chất lượng cao trong kỷ nguyên Công nghiệp 4.0.

\subsection{Yêu cầu của đề tài}
Đề tài yêu cầu khảo sát các giải pháp và sản phẩm liên quan, đồng thời phân tích, so sánh ưu - nhược điểm của từng hướng tiếp cận. Trên cơ sở đó, đề xuất hoặc lựa chọn một mô hình học sâu phù hợp với các ràng buộc về tài nguyên phần cứng. Cuối cùng, tiến hành thiết kế, hiện thực hệ thống xác thực khuôn mặt và triển khai thử nghiệm trên ít nhất một nền tảng vi xử lý nhằm đánh giá hiệu năng trong điều kiện thực tế.